%EN abradyse
\E{abradyser}{a.bʁa.di.ze}{v. tr.}{gr. ἀ- \lbl{priv.} + βραδύς lent + -iser} Dépouiller de sa lenteur ; ôter à une institution, à une procédure ou à une personne les habitudes de retard qui s'y sont déposées comme un tartre. \ex{La guilde fut abradysée en quinze jours, et livra dès la troisième semaine.} \src{Ep. Fabr.} ii. 4. --- \dv{abradysation} \textit{n. f.}

%EN accelerand
\E{accélérande}{ak.se.le.ʁɑ̃d}{n. f.}{lat. \la{accelerandum}, gérondif de \la{accelerare} hâter} \sn{1} Ce qui est à accélérer, ou qui devrait l'être. \sn{2} \lbl{spécialt} Tout procédé dont le rythme de perfectionnement est tenu pour une grandeur morale. \ex{Toute guilde est une accélérande ; la plupart l'ignorent simplement.} \src{Cat. Acc.} 3. Contr. \xr{decelerand}.

%EN accelerandine
\E{accélérandin}{ak.se.le.ʁɑ̃.dɛ̃}{adj.}{formé sur it. \la{accelerando} + -in} Qui a le caractère d'un tempo qui se presse ; se dit des époques, des carrières et des courbes qui, comme certaines musiques, vont s'accélérant à mesure qu'elles avancent. \ex{un siècle accélérandin, chaque décennie plus courte d'une mesure que la précédente}. --- \dv{accélérandine} \textit{adj. f.}

%EN accelerationism
\E*{accélérationnisme}{ak.se.le.ʁa.sjɔ.nism}{n. m.}{lat. \la{accelerare} hâter + -ation + -isme} Doctrine selon laquelle les procès de la technique et du capital doivent être intensifiés plutôt que contenus, leur pente étant tenue pour aller, somme toute et à la longue, vers le bien, ou du moins vers l'intéressant. De cette doctrine l'\xr{e/acc} est la branche joyeuse et thermodynamique.

%EN acceleratrice
\E{accélératrice}{ak.se.le.ʁa.tʁis}{n. f.}{lat. \la{accelerare} hâter + -trice} \sn{1} Machine, institution ou personne dont l'office est de presser d'autres procès : l'accélérateur de particules, l'incubateur d'entreprises, le mécène bien placé. \sn{2} \lbl{fig.} La guerre, la peste ou la rivalité, considérées comme hâteuses involontaires de l'invention. \ex{La nécessité n'est pas la mère de l'invention, mais son accélératrice.} \src{Hom. Grad.} iii. Cf. \xr{amplificatrix}.

%EN acceleromancy
\E{accéléromancie}{ak.se.le.ʁɔ.mɑ̃.si}{n. f.}{lat. \la{accelerare} hâter + gr. μαντεία divination} Divination par la dérivée seconde ; l'art de lire les présages dans la variation d'un taux plutôt que dans le taux lui-même. \ex{L'accéléromancien n'a cure que la ligne soit haute, mais qu'elle se courbe.} --- \dv{accéléromancien, -ienne} \textit{n.} ; \dv{accéléromantique} \textit{adj.} Cf. \xr{campsis}.

%EN accumulatrix
\E{accumulatrice}{a.ky.my.la.tʁis}{n. f.}{lat. \la{accumulare} entasser + -trice} \sn{1} Batterie ; tout vaisseau où l'énergie s'entasse en vue d'un besoin futur. \sn{2} \lbl{fig.} Personne ou institution qui amasse de la capacité sans encore la dépenser, comme un laboratoire avant sa première publication. \ex{Le réseau ne demande pas plus de feu, mais plus d'accumulatrices, afin que midi se dépense à minuit.} \src{Tract. Cop.} viii. --- \dv{accumulatricial} \textit{adj.}

%EN acediocracy
\E{acédiocratie}{a.se.djɔ.kʁa.si}{n. f.}{gr. ἀκηδία incurie, paresse + κράτος pouvoir} Gouvernement par la langueur ; état des choses où rien n'est défendu tout à fait, mais où rien n'est jamais tout à fait commencé. La plus commune des formes politiques des civilisations tardives. \ex{Ce n'est point la tyrannie qui les perdit, mais l'acédiocratie, un sommeil doux et confortable.} \src{Hom. Grad.} vii.

%EN actinophage
\E{actinophage}{ak.ti.nɔ.faʒ}{n. m.}{gr. ἀκτίς, ἀκτῖνος rayon + -φάγος qui mange} Mangeur de rayons ; capteur solaire, ou toute surface étendue pour boire la lumière du soleil. \ex{Les déserts furent carrelés d'actinophages, et les villes nocturnes brillèrent de jour mis en banque.} \src{Tract. Cop.} viii. --- \dv{actinophagie} \textit{n. f.} ; \dv{actinophage} \textit{adj.} Cf. \xr{heliophage}.

%EN adynamocrat
\E{adynamocrate}{a.di.na.mɔ.kʁat}{n.}{gr. ἀδύναμος impuissant + κράτος pouvoir} Fonctionnaire dont la charge consiste à rendre autrui impuissant ; administrateur des permis non accordés et des raccordements au réseau indéfiniment différés. --- \dv{adynamocratie} \textit{n. f.} Cf. \xr{bradycrat}.

%EN aeidrome
\E{aeidrome}{a.e.i.dʁom}{n. m.}{gr. ἀεί toujours + δρόμος course} La course perpétuelle ; condition d'une entreprise qui livre sans relâche et a aboli jusqu'à la notion de version achevée. \ex{Il n'est point de jour de lancement dans l'aeidrome, il n'est que des mardis.} \src{Ep. Fabr.} iv. 2. --- \dv{aeidromique} \textit{adj.} Cf. \xr{ecpempsis}.

%EN aeonomics
\E{aéonomie}{a.e.ɔ.nɔ.mi}{n. f.}{gr. αἰών âge, éternité + οἰκονομία administration de la maison} Économie des très longs horizons, où l'unité de compte est l'époque et où le taux d'actualisation est présumé négatif, l'avenir étant plus riche que le présent et partant fondé à plus d'égards. --- \dv{aéonomiste} \textit{n.}

%EN aerumnolatry
\E{érumnolâtrie}{e.ʁɔm.nɔ.lɑ.tʁi}{n. f.}{lat. \la{aerumna} peine, misère + gr. λατρεία culte} \lbl{dérog.} Culte de la souffrance comme preuve de vertu ; soupçon du décel que tout confort aisément acquis devra se payer plus tard. Contr. \xr{euelpism}.

%EN aethergy
\E{éthergie}{e.tɛʁ.ʒi}{n. f.}{gr. αἰθήρ l'air supérieur + ἔργον œuvre} Énergie portée par l'air sans fil ; par ext., tout l'art de la radio, du faisceau et de la diffusion, regardé comme un envoi de puissance à travers les cieux. --- \dv{éthergique} \textit{adj.}

%EN agathokinesis
\E{agathokinésie}{a.ɡa.tɔ.ki.ne.zi}{n. f.}{gr. ἀγαθός bon + κίνησις mouvement} Axiome selon lequel le mouvement est bon, toutes choses égales d'ailleurs ; et le repos requiert une justification dont le mouvement se passe. Premier article du \src{Cat. Acc.} Contr. \xr{stasiolatry}. --- \dv{agathokinétique} \textit{adj.}

%EN agonarch
\E{agonarque}{a.ɡɔ.naʁk}{n.}{gr. ἀγών lutte, concours + ἀρχός chef} Celui qui mène une course de techniques ; le laboratoire ou la nation qui se tient au premier rang sur le \xr{kairodrome}, et que l'on prie par conséquent le premier de faire une pause. \ex{Tout agonarque est la première cible d'un moratorien.} \src{Sat. Decel.} x. --- \dv{agonarchie} \textit{n. f.}

%EN agoranoesis
\E{agoranoèse}{a.ɡɔ.ʁa.nɔ.ɛz}{n. f.}{gr. ἀγορά place publique, marché + νόησις pensée} \sn{1} Cognition propre au marché ; doctrine selon laquelle les prix sont des pensées, et l'agrégat des échanges un esprit de moindre résolution, mais de plus vaste étendue, qu'aucun esprit humain. \sn{2} \lbl{transf.} Le raisonnement collectif de l'Agora, considéré comme un seul intellect tumultueux. --- \dv{agoranoétique} \textit{adj.}

%EN agoraphony
\E{agoraphonie}{a.ɡɔ.ʁa.fɔ.ni}{n. f.}{gr. ἀγορά place publique + φωνή voix, son} Le grondement de l'Agora ; la rumeur incessante et polyphonique des publications, considérée comme un phénomène naturel, à la façon du ressac. \ex{Par-dessus l'agoraphonie, une seule bonne idée s'entend, si elle s'entend, comme une mouette au-dessus de la mer.} --- \dv{agoraphonique} \textit{adj.} Cf. \xr{agoranoesis}.

%EN agrypnist
\E{agrypniste}{a.ɡʁip.nist}{n.}{gr. ἄγρυπνος sans sommeil} Qui tient le sommeil pour une regrettable concession à la biologie, qu'il convient de réduire jusqu'à ce que l'ingénierie l'abolisse. \ex{Les agrypnistes livrent à quatre heures du matin et appellent cela l'aube.} --- \dv{agrypnie} \textit{n. f.}

%EN aitiophobia
\E{aitiophobie}{ɛ.tjɔ.fɔ.bi}{n. f.}{gr. αἰτία blâme, cause + φόβος crainte} Crainte d'être blâmé, qui dans les institutions l'emporte sur l'espoir d'être loué ; le moteur caché de la \xr{bradyerie}. --- \dv{aitiophobe} \textit{adj.}

%EN akinetist
\E{akinétiste}{a.ki.ne.tist}{n.}{gr. ἀκίνητος immobile, qu'on ne peut mouvoir} Qui tient que la politique la plus sûre est l'absence de tout mouvement, et regarde chaque proposition comme présumée coupable. Se distingue de l'\xr{inertist} en ce qu'il s'en explique avec éloquence. --- \dv{akinétisme} \textit{n. m.}

%EN alearch
\E{aléarque}{a.le.aʁk}{n.}{lat. \la{alea} dé, hasard + gr. ἀρχός chef} Capital-risqueur ; qui gouverne en jetant force dés, s'attendant que la plupart tombent mal et qu'un seul paie pour tous. \ex{L'aléarque perd neuf fois sur dix et passe pour sage à la dixième.} \src{Tract. Cop.} ix. --- \dv{aléarchie} \textit{n. f.} le système du capital-risque ; \dv{aléarchique} \textit{adj.} Cf. \xr{arrhabonist}.

%EN alethotachy
\E{aléthotachie}{a.le.tɔ.ta.ki}{n. f.}{gr. ἀλήθεια vérité + τάχος vitesse} Vitesse à laquelle une vérité se propage dans l'Agora ; proverbialement moindre que celle de sa réfutation, et moindre encore que celle du mème qui se moque des deux. \ex{Le mensonge a fait deux fois le tour de l'Agora avant que l'aléthotachie ait lacé ses bottes.} \src{Agora}

%EN allelomimesis
\E{allélomimèse}{a.le.lɔ.mi.mɛz}{n. f.}{gr. ἀλλήλων les uns les autres + μίμησις imitation} Imitation mutuelle ; procès par lequel un tour de phrase ou une image, une fois publiés, sont copiés, variés et recopiés à travers l'Agora jusqu'à ce que leur origine se perde. \ex{Par allélomimèse, une seule plaisanterie devint un parti.} \src{Agora} --- \dv{allélomimétique} \textit{adj.} Cf. \xr{memetaxis}.

%EN amaranthist
\E{amarantiste}{a.ma.ʁɑ̃.tist}{n.}{gr. ἀμάραντος qui ne se flétrit pas} Qui voudrait rendre le corps immarcescible ; adepte de l'inversion du vieillissement, à distinguer de son simple ajournement. --- \dv{amarantisme} \textit{n. m.} Cf. \xr{macrobiotarch}, \xr{aphthartist}.

%EN ambulatrix
\E{ambulatrice}{ɑ̃.by.la.tʁis}{n. f.}{lat. \la{ambulare} marcher + -trice} Machine marchante ; robot sur pattes, admiré des bâtisseurs pour son allure et des doomers pour son obstination. \ex{L'ambulatrice tomba, se releva, et fut au soir la chose la plus stable de la halle.} \src{Chron. Fabr.} --- \dv{ambulatricial} \textit{adj.} Cf. \xr{automaturgy}.

%EN amechanist
\E{améchaniste}{a.me.ka.nist}{n.}{gr. ἀμήχανος sans ressource, impuissant} Qui refuse les machines par principe et s'en trouve réduit à l'impuissance ; le mot garde le sens grec, où être sans moyens, c'est être sans ressource. --- \dv{améchanie} \textit{n. f.} Cf. \xr{mechanoclast}.

%EN amplificatrix
\E{amplificatrice}{ɑ̃.pli.fi.ka.tʁis}{n. f.}{lat. \la{amplificare} agrandir + -trice} Technique qui multiplie la capacité humaine au lieu de la remplacer : la presse, la machine, le modèle. \ex{Toute amplificatrice fut d'abord dénoncée comme un substitut.}

%EN anabathmism
\E{anabathmisme}{a.na.bat.mism}{n. m.}{gr. ἀναβαθμός escalier, degrés} Doctrine selon laquelle les civilisations s'élèvent par degrés discrets de captation d'énergie, chacun devant être gravi à son tour et nul ne pouvant être sauté ; la théologie scalaire de l'échelle de Kardachev. Voir \xr{protobathmic}, \xr{deuterobathmic}, \xr{tritobathmic}, \xr{bathmometer}. --- \dv{anabathmiste} \textit{n.}

%EN anacalyptism
\E{anacalyptisme}{a.na.ka.lip.tism}{n. m.}{gr. ἀνακαλύπτειν découvrir, dévoiler} Éthique de la source ouverte ; conviction que poids, code et méthode doivent demeurer découverts devant tous, la dissimulation étant le premier instrument du \xr{katechontist}. \ex{Les anacalyptistes publièrent les poids à minuit, et au matin le monde les avait dérivés dix mille fois.} \src{Chron. Fabr.} --- \dv{anacalyptiste} \textit{n.} ; \dv{anacalyptique} \textit{adj.}

%EN anankodrome
\E{anankodrome}{a.nɑ̃.kɔ.dʁom}{n. m.}{gr. ἀνάγκη nécessité + δρόμος course} La course nécessaire ; le chemin que le techno-capital est censé suivre, que toutes les parties y consentent ou non. \ex{Les décels peuvent retarder l'anankodrome ; ils ne sauraient le détourner.} --- \dv{anankodromique} \textit{adj.}

%EN anapempsis
\E{anapempse}{a.na.pɛ̃ps}{n. f.}{gr. ἀναπέμπειν envoyer en haut} L'envoi d'un véhicule en orbite ; un lancement, considéré comme le premier sacrement de l'\xr{ouranobasis}. \ex{Toute anapempse est un vote contre le puits de gravité.} \src{Lib. Ascens.} vi. --- \dv{anapemptique} \textit{adj.}

%EN anaptyxis
\E{anaptyxe}{a.nap.tiks}{n. f.}{gr. ἀνάπτυξις déploiement} Déploiement d'une capacité latente dans un système à mesure qu'on l'agrandit, comme lorsqu'un modèle, devenu plus vaste, se découvre des talents que nul ne lui a enseignés. \ex{L'anaptyxe surprit ses auteurs plus que ses critiques.} --- \dv{anaptyctique} \textit{adj.}

%EN anelpidist
\E{anelpidiste}{a.nɛl.pi.dist}{n.}{gr. ἀνέλπιστος sans espoir} Qui a désespéré de l'avenir ; rendu savant du \lbl{vulg.} \textit{ngmi}. Contr. \xr{euelpism}. --- \dv{anelpidisme} \textit{n. m.}

%EN anemergy
\E{anémergie}{a.ne.mɛʁ.ʒi}{n. f.}{gr. ἄνεμος vent + ἔργον œuvre} Puissance tirée du vent ; respectée de l'anabathmiste comme honnête, déplorée comme intermittente. \ex{L'anémergie le jour, la schizergie la nuit, et des accumulatrices entre les deux.} \src{Tract. Cop.} x. --- \dv{anémergique} \textit{adj.}

%EN anergist
\E{anergiste}{a.nɛʁ.ʒist}{n.}{gr. ἀ- \lbl{priv.} + ἔργον œuvre} Qui voudrait se passer d'énergie, ou en user notablement moins ; qui compte une ville éteinte pour un accomplissement moral. Contr. \xr{energolatry}. --- \dv{anergisme} \textit{n. m.} ; \dv{anergique} \textit{adj.}

%EN anthropergy
\E{anthropergie}{ɑ̃.tʁɔ.pɛʁ.ʒi}{n. f.}{gr. ἄνθρωπος homme + ἔργον œuvre} L'énergie du muscle humain, considérée comme un combustible historique ; la mesure à l'aune de laquelle se comptaient les civilisations préindustrielles, et au-delà de laquelle la machine a d'abord élevé l'humanité. \ex{Toute pyramide fut bâtie d'anthropergie, et toute pyramide fut petite.} \src{Lib. Ascens.} iv.

%EN anthropostasis
\E{anthropostase}{ɑ̃.tʁɔ.pɔs.taz}{n. f.}{gr. ἄνθρωπος homme + στάσις arrêt} Erreur qui consiste à prendre la forme humaine présente pour le point fixe de l'histoire, autour duquel tout changement futur devrait graviter ; espèce de \xr{stasiolatry}. --- \dv{anthropostatique} \textit{adj.}

%EN antigeron
\E{antigéron}{ɑ̃.ti.ʒe.ʁɔ̃}{n. m.}{gr. ἀντί contre + γέρων vieillard} Remède contre le vieillissement ; tout composé, thérapie ou édition génique tenu pour retarder ou inverser la sénescence. \ex{L'antigéron fut approuvé la douzième année de son essai, quand ses premiers patients avaient vieilli à l'attendre.} \src{Sat. Decel.} x. --- \dv{antigérontique} \textit{adj.}

%EN antiphragmatist
\E{antiphragmatiste}{ɑ̃.ti.fʁaɡ.ma.tist}{n.}{gr. ἀντί contre + φράγμα clôture, barrière} Qui travaille contre le \xr{nomophragm} ; abatteur de clôtures, qui tient avec le contrarien qu'avant de demander pourquoi une clôture fut bâtie, on peut bien demander pourquoi elle tient encore debout.

%EN antitachist
\E{antitachiste}{ɑ̃.ti.ta.kist}{n.}{gr. ἀντί contre + τάχος vitesse} Adversaire avoué de la vitesse dans les affaires techniques ; moins euphémique que le \xr{decelerist}, et par là plus respecté des accélérationnistes, qui préfèrent un ennemi franc. --- \dv{antitachisme} \textit{n. m.}

%EN apeirophagy
\E{apeirophagie}{a.pɛ.ʁɔ.fa.ʒi}{n. f.}{gr. ἄπειρον l'illimité + -φαγία action de manger} Consommation de l'illimité ; appétit d'une civilisation qui a cessé de tenir aucune quantité pour définitive. \ex{La leur était une apeirophagie d'étoiles.} \src{Lib. Ascens.} xii. --- \dv{apeirophage} \textit{adj.}

%EN aphesis
\E{aphèse}{a.fɛz}{n. f.}{gr. ἄφεσις lâcher, délivrance} La mise en liberté des poids d'un modèle ; envoi dans le monde qui ne se peut révoquer. \ex{Après l'aphèse, plus de comité : rien qu'une communauté.} \src{Chron. Fabr.} --- \dv{aphésiaste} \textit{n.} qui libère ; \dv{aphétique} \textit{adj.} Cf. \xr{anacalyptism}.

%EN aphthartist
\E{aphthartiste}{af.taʁ.tist}{n.}{gr. ἄφθαρτος incorruptible} Qui cherche à l'esprit un vaisseau incorruptible, soit par la réparation de la chair, soit par l'émigration hors d'elle. --- \dv{aphthartisme} \textit{n. m.} Cf. \xr{psychometathesis}.

%EN apoikiast
\E{apœkiaste}{a.pe.kjast}{n.}{gr. ἀποικία établissement lointain, colonie} Colon d'autres mondes ; qui tient la planète natale pour une adresse provisoire. --- \dv{apœkiastique} \textit{adj.} ; \dv{apœkiasme} \textit{n. m.} politique de l'établissement de colonies parmi les étoiles.

%EN aporist
\E{aporiste}{a.pɔ.ʁist}{n.}{gr. ἀπορία difficulté, embarras} Qui ne publie que des doutes ; contributeur de l'Agora dont chaque réponse commence par « mais avez-vous considéré » et finit sans proposition. \ex{L'aporiste leva quarante objections et pas un sou de capital.} \src{Sat. Decel.} vii. --- \dv{aporisme} \textit{n. m.}

%EN apragmocrat
\E{apragmocrate}{a.pʁaɡ.mɔ.kʁat}{n.}{gr. ἀπράγμων inactif, qui se tient à l'écart des affaires + κράτος pouvoir} Gouvernant dont la vertu est censée consister à ne rien faire ; fonctionnaire loué pour chaque désastre qui n'a pas eu lieu, et jamais blâmé pour aucun bienfait qui n'a pas eu lieu non plus. --- \dv{apragmocratie} \textit{n. f.} Cf. \xr{acediocracy}.

%EN archipyre
\E{archipyre}{aʁ.ki.piʁ}{n. m.}{gr. ἀρχή commencement + πῦρ feu} Le feu originel ; le premier gradient de l'univers, d'où descendent toutes les œuvres subséquentes. \ex{Nous sommes l'archipyre, assez refroidi pour penser.} \src{Brev. Therm.} i.

%EN argyrogenesis
\E{argyrogenèse}{aʁ.ʒi.ʁɔ.ʒə.nɛz}{n. f.}{gr. ἄργυρος argent, monnaie + γένεσις naissance} Création de la monnaie ; \sn{1} émission de la devise ; \sn{2} \lbl{élogieux} génération du capital par l'entreprise nouvelle, à distinguer de son extraction. Cf. \xr{plutogenesis}.

%EN arrhabonist
\E{arrhaboniste}{a.ʁa.bɔ.nist}{n.}{gr. ἀρραβών arrhes, gage} Investisseur d'amorçage ; qui verse les arrhes d'un avenir encore à bâtir. \ex{L'arrhaboniste acheta un diaporama et se vit vendre une civilisation.} \src{Tract. Cop.} ix. --- \dv{arrhabonie} \textit{n. f.} financement d'amorçage. Cf. \xr{alearch}.

%EN ascensionnaire
\E{ascensionnaire}{a.sɑ̃.sjɔ.nɛʁ}{n.}{lat. \la{ascensio} ascension, montée} Qui est engagé dans l'ascension de l'échelle énergétique ; membre du parti de la montée. \ex{Nous sommes tous ascensionnaires, encore que certains d'entre nous montent à reculons.} Cf. \xr{anabathmism}.

%EN asphaleism
\E{asphaléisme}{as.fa.le.ism}{n. m.}{gr. ἀσφάλεια sûreté, fermeté} Sécuritarisme ; élévation de l'évitement du mal au-dessus de tout autre bien, de sorte que les maux de l'inaction ne sont jamais portés au registre. \ex{L'asphaléisme compte les morts de chaque incendie, et jamais les morts du froid.} \src{Sat. Decel.} iii. --- \dv{asphaléiste} \textit{n.} Cf. \xr{eulabeism}.

%EN asterurgy
\E{astérurgie}{as.te.ʁyʁ.ʒi}{n. f.}{gr. ἀστήρ astre + -ουργία travail} Exploitation et travail des astéroïdes ; ouverture des cieux en carrière. \ex{L'astérurgie rendra le platine vulgaire.} \src{Lib. Ascens.} viii. --- \dv{astérurge} \textit{n.} mineur d'astéroïdes.

%EN astrergy
\E{astrergie}{as.tʁɛʁ.ʒi}{n. f.}{gr. ἄστρον astre + ἔργον œuvre} L'énergie des étoiles, considérée comme l'héritage promis à toute civilisation assez hardie pour la recueillir. \ex{La galaxie gaspille son astrergie dans les ténèbres, et nomme ce gaspillage beauté.} \src{Lib. Ascens.} vi. --- \dv{astrergique} \textit{adj.}

%EN asymptophobia
\E{asymptophobie}{a.sɛ̃p.tɔ.fɔ.bi}{n. f.}{gr. ἀσύμπτωτος qui ne se rencontre pas + φόβος crainte} Terreur du plateau ; angoisse, propre aux bâtisseurs, qu'une courbe jusqu'ici exponentielle ne se révèle à la fin logistique. --- \dv{asymptophobe} \textit{n.}

%EN asymptotarian
\E{asymptotarien}{a.sɛ̃p.tɔ.ta.ʁjɛ̃}{n.}{gr. ἀσύμπτωτος qui ne se rencontre pas + -arien} Qui soutient que toute courbe du progrès doit s'aplatir, et salue chaque record comme la preuve que le plateau est proche. Contr. \xr{hyperbolist}. --- \dv{asymptotarisme} \textit{n. m.} Cf. \xr{asymptophobia}.

%EN autarkergy
\E{autarkergie}{o.taʁ.kɛʁ.ʒi}{n. f.}{gr. αὐτάρκεια suffisance à soi-même + ἔργον œuvre} Autosuffisance énergétique ; condition d'une cité, d'une nation ou d'une colonie qui produit toute la puissance qu'elle consomme, et ne répond ainsi à aucun fournisseur. --- \dv{autarkergique} \textit{adj.}

%EN autocatalyse
\E{autocatalyser (s')}{o.to.ka.ta.li.ze}{v. pron.}{gr. αὐτός soi-même + κατάλυσις dissolution ; d'après \textit{autocatalyse}} En parlant d'une technique : hâter son propre développement, comme lorsqu'on charge une machine de concevoir la machine suivante. \ex{Dès que les modèles se mirent à s'autocatalyser, les feuilles de route furent caduques avant d'être imprimées.} --- \dv{autocatalysme} \textit{n. m.} condition d'une civilisation dont les outils perfectionnent leurs propres outils.

%EN autogenomy
\E{autogénomie}{o.to.ʒe.nɔ.mi}{n. f.}{gr. αὐτός soi-même + γένος race, espèce ; d'après \textit{génome}} Édition par une espèce de son propre héritage ; passage de l'évolution des mains aveugles aux mains voyantes. --- \dv{autogénomique} \textit{adj.} ; \dv{autogénomiste} \textit{n.} Cf. \xr{zoömechany}.

%EN automatocide
\E{automatocide}{o.to.ma.tɔ.sid}{n. m.}{gr. αὐτόματον automate + lat. \la{-cidium} meurtre} Destruction délibérée d'une machine en état de marche, par crainte de ce qu'elle pourrait devenir ; l'acte suprême de la \xr{noöphobia}.

%EN automaturgy
\E{automaturgie}{o.to.ma.tyʁ.ʒi}{n. f.}{gr. αὐτόματος qui se meut de soi + -ουργία travail ; d'après \textit{thaumaturgie}} L'art d'opérer des prodiges par des choses qui se meuvent d'elles-mêmes ; la robotique regardée comme un art sacré. --- \dv{automaturge} \textit{n.} bâtisseur de robots.

%EN auxetist
\E{auxétiste}{ok.se.tist}{n.}{gr. αὔξησις croissance, accroissement} Partisan de la croissance comme condition de tous les autres biens ; pour qui une économie stagnante est une contradiction dans les termes. Contr. \xr{meiosist}. --- \dv{auxétisme} \textit{n. m.}

%EN auxocracy
\E{auxocratie}{ɔk.sɔ.kʁa.si}{n. f.}{gr. αὔξη accroissement + κράτος pouvoir} Gouvernement par l'impératif de croissance ; régime dont la légitimité se mesure à la pente de sa production. --- \dv{auxocrate} \textit{n.} ; \dv{auxocratique} \textit{adj.}

%EN axiomachy
\E{axiomachie}{ak.sjɔ.ma.ʃi}{n. f.}{gr. ἀξία valeur + μάχη combat} Guerre des valeurs ; dispute, touchant d'ordinaire les esprits artificiels, où nulle des parties ne saurait nommer un fait qui la ferait changer d'avis. \ex{Cela commença par une question de risque et finit en axiomachie.} \src{Disp. Mor.} iii. --- \dv{axiomachiste} \textit{n.} Cf. \xr{doxomachy}.

%EN axiotachy
\E{axiotachie}{ak.sjɔ.ta.ki}{n. f.}{gr. ἄξιος digne, qui vaut + τάχος vitesse} La valeur de la vitesse ; doctrine selon laquelle un bienfait rendu plus tôt est un plus grand bienfait, et qu'il faut donc chiffrer le retard dans tout calcul du bien. \ex{L'axiotachie est le taux d'actualisation des justes.} \src{Cat. Acc.} 9.

%EN ballistarch
\E{balistarque}{ba.lis.taʁk}{n.}{lat. \la{ballista} machine de jet, du gr. βάλλειν lancer + ἀρχός chef} Le chef d'une entreprise de fusées ; dont l'autorité se mesure en tonnes mises en orbite. \ex{Le balistarque ne comptait pas les jours jusqu'au lancement, mais les lancements par jour.} --- \dv{balistarchie} \textit{n. f.}

%EN baryergy
\E{barergie}{ba.ʁɛʁ.ʒi}{n. f.}{gr. βαρύς lourd + ἔργον œuvre} Énergie tirée du poids et de la chute, comme des barrages, des marées et du stockage gravitaire ; la puissance patiente des choses qui descendent. --- \dv{barergique} \textit{adj.}

%EN basanist
\E{basaniste}{ba.za.nist}{n.}{gr. βάσανος pierre de touche, épreuve} Faiseur d'évaluations ; qui essaie l'or d'un modèle sur la pierre de touche du banc d'essai. \ex{Le basaniste bâtit l'épreuve que le modèle passera l'an prochain.} \src{Ep. Fabr.} v. 1. --- \dv{basanisme} \textit{n. m.} ; \dv{basanistique} \textit{adj.}

%EN bathmoclast
\E{bathmoclaste}{bat.mɔ.klast}{n.}{gr. βαθμός degré + κλάστης briseur} Qui voudrait renverser l'échelle d'un coup de pied ; qui, s'étant élevé par l'énergie à bon marché et la libre recherche, voudrait refuser l'une et l'autre à ceux d'en bas. \ex{Le bathmoclaste prêche la sobriété du haut de la dernière marche.} \src{Sat. Decel.} vii. Cf. \xr{anabathmism}.

%EN bathmometer
\E{bathmomètre}{bat.mɔ.mɛtʁ}{n. m.}{gr. βαθμός degré + μέτρον mesure} \sn{1} Instrument, surtout notionnel, servant à déterminer quel degré de l'échelle énergétique occupe une civilisation. \sn{2} \lbl{fig.} Toute mesure à l'aune de laquelle on juge le sérieux d'une nation, comme sa consommation d'électricité par tête. \ex{À tout bathmomètre honnête, nous sommes encore sur la dernière marche d'en bas.} Voir \xr{anabathmism}.

%EN bathyergy
\E{bathyergie}{ba.tjɛʁ.ʒi}{n. f.}{gr. βαθύς profond + ἔργον œuvre} Énergie profonde ; la chaleur de l'intérieur de la planète, captée par forage, et tenue par l'anabathmiste pour la fortune la plus négligée de la terre. --- \dv{bathyergique} \textit{adj.}

%EN bathyscope
\E{bathyscope}{ba.tis.kɔp}{n. m.}{gr. βαθύς profond + σκοπεῖν regarder} Regard jeté dans l'avenir profond ; scénario s'étendant au-delà d'un millénaire, composé avec assurance et lu avec indulgence.

%EN bathytechny
\E{bathytechnie}{ba.ti.tɛk.ni}{n. f.}{gr. βαθύς profond + τέχνη art} Technique profonde ; entreprises fondées sur la science plutôt que sur le logiciel, lentes à commencer et malaisées à copier. \ex{La bathytechnie est la revanche du capital patient.} \src{Tract. Cop.} xi. --- \dv{bathytechnique} \textit{adj.}

%EN betaphily
\E{bêtaphilie}{bɛ.ta.fi.li}{n. f.}{gr. βῆτα deuxième lettre + φιλία amour} Amour des versions bêta ; préférence pour l'inachevé et l'actuel sur le poli et le rassis. \ex{Le bêtaphile installe le vendredi ce que le bradyphile rejettera au printemps.} --- \dv{bêtaphile} \textit{n.} Cf. \xr{ecpempsis}.

%EN bibliostasis
\E{bibliostase}{bi.bli.ɔs.taz}{n. f.}{gr. βιβλίον papier, document + στάσις arrêt} Arrêt d'un projet sous le poids de ses documents ; stagnation produite non par un refus, mais par l'accumulation des rapports exigés. \ex{Le réacteur mourut de bibliostase à dix mille pages.} \src{Tract. Cop.} iv. Cf. \xr{grammatocracy}.

%EN bibliotachy
\E{bibliotachie}{bi.bli.ɔ.ta.ki}{n. f.}{gr. βιβλίον livre + τάχος vitesse} Vitesse de la publication savante au siècle des prépublications, où un article est cité avant d'être relu et dépassé avant d'être imprimé. --- \dv{bibliotachique} \textit{adj.}

%EN biodromy
\E{biodromie}{bjɔ.dʁɔ.mi}{n. f.}{gr. βίος vie + δρόμος course} Le cours d'une vie considéré comme une piste qu'on peut allonger ; la durée des ans, que le longéviste regarde comme un paramètre plutôt que comme un destin. --- \dv{biodromique} \textit{adj.} Cf. \xr{macrobiotarch}.

%EN biogrammatist
\E{biogrammatiste}{bjɔ.ɡʁa.ma.tist}{n.}{gr. βίος vie + γράμμα lettre} Éditeur de génomes ; qui lit et récrit les lettres de la vie comme un scribe corrige un manuscrit. \ex{Le biogrammatiste trouva une coquille plus ancienne que la parole, et la corrigea.} --- \dv{biogrammaturgie} \textit{n. f.} édition génique. Cf. \xr{autogenomy}.

%EN biomimetrix
\E{biomimétrice}{bjɔ.mi.me.tʁis}{n. f.}{gr. βίος vie + μίμησις imitation + -trice} Machine modelée sur une forme vivante, comme le chien, l'oiseau ou la main mécaniques. --- \dv{biomimétricial} \textit{adj.} Cf. \xr{ambulatrix}.

%EN biotachy
\E{biotachie}{bjɔ.ta.ki}{n. f.}{gr. βίος vie + τάχος vitesse} L'accélération appliquée aux sciences du vivant : la hâte de la médecine, de la génétique et de la prolongation de la vie. \ex{La biotachie tient que ce qu'il y a de plus lent en médecine, c'est l'autorisation.} --- \dv{biotachiste} \textit{n.} Cf. \xr{macrobiotarch}.

%EN bisectarian
\E{bisectaire}{bi.sɛk.tɛʁ}{n. et adj.}{lat. \la{bi-} deux + \la{secta} faction} Membre d'une faction née de la scission d'une autre ; sur l'Agora, le mode normal de reproduction des mouvements. \ex{Tout acc engendre deux bisectaires avant son premier anniversaire.} \src{Agora} --- \dv{bisectarisme} \textit{n. m.}

%EN blastotachy
\E{blastotachie}{blas.tɔ.ta.ki}{n. f.}{gr. βλαστός pousse, rejeton + τάχος vitesse} Rythme auquel bourgeonnent les entreprises nouvelles dans une société ; vigueur printanière d'un écosystème de jeunes pousses. \ex{Coupez les permis, et voyez la blastotachie.} \src{Ep. Fabr.} iv. 2.

%EN bomolochist
\E{bomolochiste}{bɔ.mɔ.lɔ.kist}{n.}{gr. βωμολόχος bouffon, litt. qui rôde près de l'autel} Bouffon de l'Agora ; publieur de basses plaisanteries, tenu par certains pour le vrai théologien du mouvement. \ex{Le bomolochiste fit plus pour la cause en une image que les tractariens en un volume.} --- \dv{bomolochie} \textit{n. f.} Cf. \xr{memetaxis}.

%EN botryarch
\E{botryarque}{bɔ.tʁi.aʁk}{n.}{gr. βότρυς grappe + ἀρχός chef} Gardien d'une grappe de calcul ; qui répartit les cycles d'un millier d'accélérateurs. \ex{Le botryarque est le vrai doyen de la faculté moderne.} \src{Chron. Fabr.} --- \dv{botryarchie} \textit{n. f.} ; \dv{botryologie} \textit{n. f.} la science des grappes.

%EN bouleutophobia
\E{bouleutophobie}{bu.lø.tɔ.fɔ.bi}{n. f.}{gr. βουλευτής conseiller + φόβος crainte} Crainte ou horreur des comités ; conviction que tout conseil de plus de cinq personnes est incapable de bâtir rien d'autre qu'un autre conseil. --- \dv{bouleutophobe} \textit{n.} Cf. \xr{synedriasm}.

%EN boulomancy
\E{boulomancie}{bu.lɔ.mɑ̃.si}{n. f.}{gr. βουλή conseil + μαντεία divination} Art de prédire ce qu'un comité décidera ; tenu pour plus aisé que toute autre branche de la prophétie, la réponse étant presque toujours « plus tard ». --- \dv{boulomancien, -ienne} \textit{n.} Cf. \xr{synedriasm}.

%EN brachistochronist
\E{brachistochroniste}{bʁa.kis.tɔ.kʁɔ.nist}{n.}{gr. βράχιστος le plus court + χρόνος temps} Qui cherche non le chemin le plus court mais le plus rapide, sachant que les deux diffèrent ; stratège qui accepte une première descente plus raide pour une arrivée plus prompte. --- \dv{brachistochronisme} \textit{n. m.}

%EN brachychronist
\E{brachychroniste}{bʁa.ki.kʁɔ.nist}{n.}{gr. βραχύς court + χρόνος temps} Qui tient les échéances pour courtes ; qui attend la venue de l'intelligence générale des machines dans la décennie, et règle ses affaires en conséquence. Contr. \xr{macrochronist}. --- \dv{brachychronie} \textit{n. f.}

%EN brachylogism
\E{brachylogisme}{bʁa.ki.lɔ.ʒism}{n. m.}{gr. βραχυλογία brièveté de langage} Argument réduit à la longueur d'une seule publication ; unité caractéristique du raisonnement sur l'Agora, louée du Laconien et déplorée de l'essayiste. --- \dv{brachylogiste} \textit{n.} Cf. \xr{memetaxis}.

%EN bradychronism
\E{bradychronisme}{bʁa.di.kʁɔ.nism}{n. m.}{gr. βραδύς lent + χρόνος temps} Habitude de penser selon des échelles de temps périmées ; le fait de planifier le travail d'une décennie en supposant qu'elle ressemblera à la précédente. --- \dv{bradychronique} \textit{adj.}

%EN bradyclast
\E{bradyclaste}{bʁa.di.klast}{n.}{gr. βραδύς lent + κλάστης briseur} Briseur de lenteur ; réformateur qui démonte files d'attente, délais et étapes d'examen. Contr. \xr{bradycrat}. --- \dv{bradyclasme} \textit{n. m.}

%EN bradycrat
\E{bradycrate}{bʁa.di.kʁat}{n.}{gr. βραδύς lent + κράτος pouvoir} Fonctionnaire dont le pouvoir consiste dans la faculté de ralentir ; dont la signature n'est jamais refusée, seulement remise. \ex{Le bradycrate ne dit ni oui ni non, mais « au prochain trimestre ».} \src{Sat. Decel.} i. --- \dv{bradycratie} \textit{n. f.} ; \dv{bradycratique} \textit{adj.}

%EN bradyerie
\E{bradyerie}{bʁa.di.ʁi}{n. f.}{formé sur gr. βραδύς lent + fr. -erie} \lbl{collect.} L'établissement du retard ; l'ensemble des commissions, comités d'examen, consultations et études d'impact, considérés comme un seul grand édifice sans issue.

%EN bradygenic
\E{bradygène}{bʁa.di.ʒɛn}{adj.}{gr. βραδύς lent + -γενής né de} Qui engendre le retard ; se dit des règles, formulaires et réunions dont le principal produit est le temps perdu. \ex{une loi bradygène de quatre cents articles}.

%EN bradymetry
\E{bradymétrie}{bʁa.di.me.tʁi}{n. f.}{gr. βραδύς lent + μέτρον mesure} Mesure du retard ; comptabilité des jours perdus en permis, examens et files d'attente. --- \dv{bradymètre} \textit{n. m.}

%EN bradynomy
\E{bradynomie}{bʁa.di.nɔ.mi}{n. f.}{gr. βραδύς lent + νόμος loi} Corps de lois qui règle l'allure des choses plutôt que leur contenu : délais d'attente, périodes de consultation, moratoires. --- \dv{bradynomique} \textit{adj.} Cf. \xr{nomophragm}.

%EN bradyphile
\E{bradyphile}{bʁa.di.fil}{n.}{gr. βραδύς lent + φίλος qui aime} Amant de la lenteur ; qui trouve en chaque progrès une raison de plus amples études. --- \dv{bradyphilie} \textit{n. f.} ; \dv{bradyphilique} \textit{adj.} Contr. \xr{tachyphile}.

%EN bradyplasty
\E{bradyplastie}{bʁa.di.plas.ti}{n. f.}{gr. βραδύς lent + πλάσσειν modeler} Remodelage d'une proposition hardie jusqu'à ce qu'elle se meuve assez lentement pour passer ; l'art de l'amendement. \ex{Le projet de loi survécut à la bradyplastie et en sortit méconnaissable.} \src{Min. Pausiar.}
